\documentclass[12pt,a4paper]{article}

\usepackage[T1]{fontenc}
\usepackage[utf8]{inputenc}
\usepackage[french]{babel}
\usepackage[a4paper,margin=2.5cm]{geometry}
\usepackage{lmodern}
\usepackage{setspace}
\usepackage{enumitem}
\usepackage{tabularx}
\usepackage{longtable}
\usepackage{array}
\usepackage{booktabs}
\usepackage{xcolor}
\usepackage{hyperref}
\usepackage{fancyhdr}
\usepackage{titlesec}

\hypersetup{
  colorlinks=true,
  linkcolor=blue!50!black,
  urlcolor=blue!50!black,
  pdftitle={Cahier des charges - Plateforme Numerique de Pharmacopée Traditionnelle Camerounaise},
  pdfauthor={Projet RICTD - Universite de Yaounde I}
}

\setstretch{1.15}
\setlength{\parindent}{0pt}
\setlength{\parskip}{6pt}
\pagestyle{fancy}
\fancyhf{}
\lhead{Cahier des charges}
\rhead{Pharmacopée Camerounaise}
\cfoot{\thepage}

\titleformat{\section}{\large\bfseries\color{blue!50!black}}{\thesection}{0.7em}{}
\titleformat{\subsection}{\normalsize\bfseries}{\thesubsection}{0.7em}{}

\begin{document}

\begin{titlepage}
  \centering
  {\Large \textbf{Université de Yaoundé I}}\\[0.4cm]
  {\large RICTD -- Numérique Universitaire}\\[2cm]

  {\Huge \textbf{Cahier des charges}}\\[0.4cm]
  {\Large Plateforme Numérique de Pharmacopée Traditionnelle Camerounaise}\\[1.4cm]

  \begin{tabular}{>{\bfseries}l l}
    Projet & Pharmacopée Camerounaise \\
    Version & 1.0 \\
    Date & \today \\
    Auteur(s) & [À compléter] \\
    Encadreur(s) & [À compléter]
  \end{tabular}

  \vfill
  \textit{Document de référence pour le cadrage fonctionnel, technique et organisationnel du projet}
\end{titlepage}

\tableofcontents
\newpage

\section{Contexte et justification}
La médecine traditionnelle occupe une place importante dans les pratiques de santé au Cameroun. Le projet vise à centraliser et valoriser les connaissances liées aux plantes médicinales, tout en proposant une plateforme numérique accessible pour les patients, chercheurs, tradipraticiens et administrateurs.

La plateforme permet notamment :
\begin{itemize}[leftmargin=1.2cm]
  \item la consultation de fiches plantes et remèdes ;
  \item la recherche par mots-clés (plante, maladie, symptôme) ;
  \item la localisation de tradipraticiens ;
  \item la gestion des témoignages avec modération ;
  \item un espace d'administration et un espace tradipraticien.
\end{itemize}

\section{Objectifs du projet}
\subsection{Objectif général}
Concevoir une application web fiable et ergonomique pour diffuser des informations de pharmacopée traditionnelle camerounaise et faciliter l'accès aux ressources de santé traditionnelle.

\subsection{Objectifs spécifiques}
\begin{itemize}[leftmargin=1.2cm]
  \item Référencer les plantes médicinales, maladies associées et usages traditionnels.
  \item Mettre à disposition une recherche simple et rapide.
  \item Offrir une cartographie des tradipraticiens validés.
  \item Assurer un contrôle qualité via validation des profils et modération des témoignages.
  \item Produire des statistiques d'usage et un export CSV pour exploitation.
\end{itemize}

\section{Périmètre}
\subsection{Inclus dans le périmètre}
\begin{itemize}[leftmargin=1.2cm]
  \item Frontend web (HTML/CSS/JS) : accueil, recherche, fiches plantes, carte, authentification, tableaux de bord.
  \item API backend PHP organisée par ressources (\texttt{/api/auth}, \texttt{/api/plantes}, \texttt{/api/tradipraticiens}, etc.).
  \item Base de données MySQL/MariaDB (schéma relationnel complet).
  \item Gestion des rôles : patient, tradipraticien, chercheur, administrateur.
\end{itemize}

\subsection{Hors périmètre (version actuelle)}
\begin{itemize}[leftmargin=1.2cm]
  \item Application mobile native.
  \item Téléconsultation en temps réel.
  \item Intégration avec un SI hospitalier externe.
\end{itemize}

\section{Parties prenantes et utilisateurs}
\begin{longtable}{p{3.5cm}p{11cm}}
\toprule
\textbf{Acteur} & \textbf{Responsabilités / Besoins} \\
\midrule
Patient & Consulter les plantes/remèdes, rechercher des solutions, laisser des témoignages. \\
Tradipraticien & Gérer son profil (géolocalisation, disponibilités), présenter ses spécialités, recevoir des avis. \\
Chercheur & Exploiter les données documentées (plantes, ethnies, maladies). \\
Administrateur & Valider les comptes tradipraticiens, modérer les témoignages, gérer les plantes, suivre les statistiques. \\
Équipe projet & Développer, maintenir, sécuriser et faire évoluer la plateforme. \\
\bottomrule
\end{longtable}

\section{Exigences fonctionnelles}
\subsection{Gestion des comptes et authentification}
\begin{itemize}[leftmargin=1.2cm]
  \item Inscription et connexion utilisateur.
  \item Gestion des rôles et contrôle d'accès par profil.
  \item Consultation de session active (\texttt{/api/auth/me.php}) et déconnexion.
\end{itemize}

\subsection{Gestion des plantes et remèdes}
\begin{itemize}[leftmargin=1.2cm]
  \item CRUD des plantes (création, lecture, mise à jour, suppression).
  \item Association des plantes aux maladies (remèdes, dosage, précautions).
  \item Affichage des fiches détaillées avec image et description.
\end{itemize}

\subsection{Recherche et assistance}
\begin{itemize}[leftmargin=1.2cm]
  \item Recherche textuelle par nom de plante, maladie ou mots associés.
  \item Assistant conversationnel pour guider l'utilisateur (\texttt{/api/chatbot/analyse.php}).
\end{itemize}

\subsection{Tradipraticiens et géolocalisation}
\begin{itemize}[leftmargin=1.2cm]
  \item Création/mise à jour du profil tradipraticien.
  \item Système de validation administrative (en attente, validé, refusé).
  \item Carte interactive des tradipraticiens validés.
\end{itemize}

\subsection{Témoignages et modération}
\begin{itemize}[leftmargin=1.2cm]
  \item Soumission de témoignages et notes.
  \item Modération des témoignages avant publication.
\end{itemize}

\subsection{Administration et pilotage}
\begin{itemize}[leftmargin=1.2cm]
  \item Tableau de bord administrateur.
  \item Statistiques globales : utilisateurs, plantes, tradipraticiens validés, témoignages approuvés.
  \item Export des données en CSV.
\end{itemize}

\section{Exigences non fonctionnelles}
\begin{itemize}[leftmargin=1.2cm]
  \item \textbf{Performance} : réponse API fluide pour les cas d'usage courants.
  \item \textbf{Sécurité} : mots de passe hachés, contrôle d'accès par rôle, validation des entrées.
  \item \textbf{Disponibilité} : exécution locale robuste sous XAMPP (Apache + MySQL).
  \item \textbf{Maintenabilité} : architecture claire (\texttt{public/}, \texttt{api/}, \texttt{includes/}, \texttt{database/}).
  \item \textbf{Ergonomie} : interface web simple, lisible et responsive.
  \item \textbf{Compatibilité} : navigateurs web modernes.
\end{itemize}

\section{Architecture technique}
\subsection{Stack technologique}
\begin{itemize}[leftmargin=1.2cm]
  \item Frontend : HTML5, CSS3, JavaScript vanilla.
  \item Backend : PHP (API REST).
  \item Base de données : MySQL / MariaDB.
  \item Cartographie : Leaflet + OpenStreetMap.
  \item Environnement local : XAMPP.
\end{itemize}

\subsection{Organisation des composants}
\begin{longtable}{p{3.5cm}p{11cm}}
\toprule
\textbf{Dossier} & \textbf{Description} \\
\midrule
\texttt{public/} & Interfaces utilisateur (pages HTML, CSS, JS, images). \\
\texttt{api/} & Endpoints backend par domaine fonctionnel. \\
\texttt{includes/} & Fonctions communes, connexion BDD, vérification d'authentification. \\
\texttt{database/} & Schéma SQL et données de démonstration. \\
\bottomrule
\end{longtable}

\section{Modèle de données (synthèse)}
Entités principales :
\begin{itemize}[leftmargin=1.2cm]
  \item \texttt{UTILISATEUR}
  \item \texttt{TRADIPRATICIEN}, \texttt{CHERCHEUR}
  \item \texttt{PLANTE}, \texttt{MALADIE}, \texttt{REMEDE}
  \item \texttt{ETHNIE}, \texttt{REGION}, \texttt{PLANTE\_ETHNIE}
  \item \texttt{TEMOIGNAGE}, \texttt{SPECIALITE}, \texttt{TRADIPRATICIEN\_SPECIALITE}
\end{itemize}

Le schéma relationnel est fourni dans \texttt{database/schema.sql}.

\section{Règles de gestion clés}
\begin{enumerate}[leftmargin=1.2cm]
  \item Un tradipraticien ne peut être visible sur la carte que s'il est validé.
  \item Un témoignage n'est affiché publiquement qu'après approbation.
  \item Les mots de passe ne sont jamais stockés en clair.
  \item Chaque compte utilisateur possède un rôle définissant ses droits.
  \item Les relations plantes-maladies sont documentées avec préparation, dosage et précautions.
\end{enumerate}

\section{Contraintes et hypothèses}
\begin{itemize}[leftmargin=1.2cm]
  \item Hébergement local initial sous XAMPP.
  \item Connexion BDD paramétrée dans \texttt{includes/db.php}.
  \item Langue principale de l'application : français.
  \item Les données de démonstration sont présentes pour accélérer les tests.
\end{itemize}

\section{Livrables attendus}
\begin{itemize}[leftmargin=1.2cm]
  \item Code source complet du projet.
  \item Base de données (script SQL).
  \item Documentation technique et utilisateur.
  \item Cahier des charges validé (présent document).
\end{itemize}

\section{Recette et critères d'acceptation}
\subsection{Exemples de scénarios de recette}
\begin{itemize}[leftmargin=1.2cm]
  \item Un utilisateur s'inscrit, se connecte et accède à son espace selon son rôle.
  \item L'administrateur valide un tradipraticien, qui devient visible sur la carte.
  \item L'administrateur modère un témoignage, qui apparaît ensuite dans la plateforme.
  \item Le module de recherche retourne des plantes pertinentes.
  \item Le tableau de bord admin affiche les statistiques et permet l'export CSV.
\end{itemize}

\subsection{Conditions de validation}
\begin{itemize}[leftmargin=1.2cm]
  \item Fonctionnalités critiques exécutables sans erreur bloquante.
  \item Données cohérentes en base après opérations CRUD.
  \item Respect des droits d'accès selon les rôles.
\end{itemize}

\section{Planning prévisionnel}
\begin{longtable}{p{2.8cm}p{11.7cm}}
\toprule
\textbf{Phase} & \textbf{Contenu} \\
\midrule
Cadrage & Analyse des besoins, définition des acteurs et du périmètre. \\
Conception & Modélisation de la base et architecture applicative. \\
Développement & Implémentation frontend, API backend et intégration BDD. \\
Tests/Recette & Vérification fonctionnelle, correction des anomalies. \\
Déploiement & Mise en production ou packaging de livraison. \\
\bottomrule
\end{longtable}

\section{Risques et mesures de mitigation}
\begin{longtable}{p{5cm}p{9.5cm}}
\toprule
\textbf{Risque} & \textbf{Mesure proposée} \\
\midrule
Qualité des données collectées & Mettre en place des validations et une modération stricte. \\
Accès non autorisé & Renforcer contrôle d'accès, sessions et vérifications serveur. \\
Montée en charge future & Prévoir optimisation SQL, indexation, pagination et cache léger. \\
Évolution des besoins & Documenter les changements et versionner les livrables. \\
\bottomrule
\end{longtable}

\section{Conclusion}
Ce cahier des charges formalise les exigences essentielles de la Plateforme Numérique de Pharmacopée Traditionnelle Camerounaise. Il sert de référence commune pour le développement, la validation et l'évolution du projet.

\vspace{0.8cm}
\textbf{Validation :}
\begin{itemize}[leftmargin=1.2cm]
  \item Responsable projet : \dotfill
  \item Encadreur : \dotfill
  \item Date : \dotfill
\end{itemize}

\end{document}
